\section{Invariant sections from a curvature-null base}
\label{inv:section}

The ordinary conversion theorem does not control characters.  We now
construct the base from invariant ratios and retain an invariant
rational section throughout the return step.  The cohomological input
comes from the natural anticanonical index.  Here the natural
linearization on $-K_X=\det T_X$ is the action induced by the
differential of the action on $X$.  The subsequent determinant,
base, and metric constructions are local to the present argument.

\begin{proposition}\label{inv:forms}
Let $X$ be a compact K\"ahler manifold, let a compact real torus $T_c$
act holomorphically, and give $L=-K_X$ its natural linearization.
Suppose $L$ has a smooth semipositive metric and
\[
 I_{T_c}(0):=\sum_q(-1)^q\dim H^q(X,\mathcal O_X)^{T_c}\ne0.
\]
There are a fixed $p\ge0$ and unbounded positive integers $m$ with
$H^0(X,\Omega_X^p\otimes(m+1)L)^{T_c}\ne0$.
\end{proposition}
\begin{proof}
The invariant anticanonical index theorem \cite[Theorem~1.1]{CompanionI}
applies to the compact
complex manifold with its natural anticanonical linearization.  It
gives an integer $a>0$ and a polynomial $P$ such that
$I_{T_c}(m)=P(m)$ for every positive multiple of $a$, with the actual
value $P(0)=I_{T_c}(0)$.  In particular $P$ is not the zero polynomial,
so one degree $q$ supplies nonzero invariant cohomology at unbounded
positive divisible exponents.

Average a K\"ahler form and the logarithmic weights of the given
metric over $T_c$.  Hard Lefschetz with semipositive coefficients
\cite[Theorem~0.1]{DPS2001} gives the equivariant surjection
\[
 H^0(X,\Omega_X^{n-q}\otimes(m+1)L)
 \longrightarrow H^q(X,K_X+(m+1)L)=H^q(X,mL).
\]
The coefficient is smooth, so its multiplier ideal is trivial.
Compact averaging makes the map surjective on invariant vectors.
Set $p=n-q$.  The shift by one is necessary because the canonical
factor cancels one copy of the naturally linearized $L$.
\end{proof}

\begin{theorem}\label{inv:nonvanishing}
Let $X$ be smooth connected projective with
$\chi(X,\mathcal O_X)\ne0$.  Let an algebraic torus $T$ act on $X$,
and suppose $L=-K_X$ is smoothly semipositive.  For the natural
anticanonical linearization,
$H^0(X,mL)^T\ne0$ for some positive integer $m$.
\end{theorem}

A connected torus acts trivially on ordinary cohomology by isotopy and
hence on $H^q(X,\mathcal O_X)$ by Hodge theory.  Thus
$\chi(X,\mathcal O_X)$ is exactly the zero value in
Proposition~\ref{inv:forms}.  The point case is immediate.  For
$n=\dim X>0$, average the metric over the maximal compact torus and
write its curvature as $\theta\ge0$, normalized to represent $c_1(L)$.
Let $\nu$ be its maximum rank and let $\omega_0$ be a K\"ahler form.
The invariant forms supplied by Proposition~\ref{inv:forms} give the
answer immediately if their degree is zero.  Otherwise determinant
extraction as in Lemma~\ref{conv:determinant} gives a linearized line
$A$, a nonzero map $A\to\bigwedge^e\Omega_X^p$, and invariant sections
\begin{equation}\label{inv:determinant-sequence}
 0\ne\sigma_j\in H^0(X,A+m_jL)^T,
 \qquad 0<m_1<m_2<\cdots.
\end{equation}
We will use an intersection consequence of the determinant map, rather
than replace this construction by the ordinary conversion theorem.

\begin{lemma}\label{inv:null-slope}
If $\nu<n$, the determinant line just constructed satisfies
\[
 c_1(A)\cdot[\theta]^\nu[\omega_0]^{n-1-\nu}\le0.
\]
\end{lemma}
\begin{proof}
For each $t>0$, Yau's prescribed-Ricci theorem \cite{Yau1978} gives
a K\"ahler metric in $[\omega_0+t\theta]$ with Ricci form
$2\pi\theta$.  The induced mean curvature on
$\bigwedge^e\Omega_X^p$ is negative semidefinite.  If $A$ had
positive degree in this K\"ahler class, solving the scalar Laplace
equation would give it a metric with constant positive curvature
trace.  The induced mean curvature on
$\operatorname{Hom}(A,\bigwedge^e\Omega_X^p)$ would then be negative
definite, contradicting the nonzero holomorphic map by the integrated
Bochner formula.  Hence
$c_1(A)\cdot[\omega_0+t\theta]^{n-1}\le0$ for every $t>0$.
Because $\theta^{\nu+1}=0$, the leading coefficient as $t\to\infty$
is the asserted nonpositive intersection, up to a positive binomial
factor.
\end{proof}

\subsection{The invariant ratio field}

Put $B_j=A+m_jL$.  Form the subfield generated by ratios of invariant
sections of the same line bundle $\sum_j k_jB_j$, where $k_j\ge0$
are integers of finite support.  Let $F$ be its relative algebraic
closure in $\mathbb C(X)$.  Both fields are finitely generated:
choose a transcendence basis for the smaller field; the algebraic
closure of the associated purely transcendental field in the finitely
generated extension $\mathbb C(X)$ is finite, and the intervening
algebraic fields are finite as well.  The field $F$ is fixed by $T$.
Indeed an element algebraic over the invariant subfield has finite
orbit among the roots of a polynomial, and connectedness makes this
orbit a point.

Finitely many ratios can be put over a common denominator by
multiplication of sections.  Thus one bundle $B=\sum k_jB_j$ and
invariant sections of $B$ define an image $Y_0$ with the smaller
function field.  Normalize in $F$, resolve the base, and obtain the
invariant rational map $X\dashrightarrow Y$, with
$\mathbb C(Y)=F$ and geometrically integral generic fiber.
Let $H_0$ be the pullback of the image polarization; it is big and
base point free.  On a smooth equivariant graph resolution
\[
 X\xleftarrow{\ \pi\ }Z\xrightarrow{\ g\ }Y
\]
there is a nonzero invariant section of
$\pi^*B-g^*H_0$: divide the pulled-back defining sections by the
homogeneous coordinates on the image.  Hence for a chosen ample $H$
on $Y$ there is an integer $b>0$ and an invariant section of
$b\pi^*B-g^*H$.

If $F=\mathbb C$, consider
\begin{equation}\label{inv:ratio-identity}
 r_j=\frac{\sigma_j^{m_2-m_1}\sigma_1^{m_j-m_2}}
                 {\sigma_2^{m_j-m_1}},\qquad j>2.
\end{equation}
Its numerator and denominator are invariant sections of the same
multiple of $B_2$, so $r_j$ is constant.  Divisor comparison and
$m_j\to\infty$ show that
$\operatorname{div}(\sigma_2/\sigma_1)\ge0$.  This is the desired
invariant section of $(m_2-m_1)L$.  We may therefore suppose
$\dim Y>0$.

There is a horizontal strictness patch on $Z$.  If $\nu=n$, use any
maximum-rank point away from the exceptional and critical loci.
If $\nu<n$, Lemma~\ref{inv:null-slope} gives
$c_1(bB)\cdot[\theta]^\nu[\omega_0]^{n-1-\nu}\le0$.
Since $b\pi^*B-g^*H$ is effective, and all the forms involved are
semipositive, it follows that
\[
 \int_Z g^*\omega_H\wedge
           \pi^*(\theta^\nu\wedge\omega_0^{n-1-\nu})=0.
\]
At a maximum-rank point where $\pi$ is an isomorphism, nonnegative
linear algebra implies $\ker\pi^*\theta\subset\ker dg$.
Choose such a point over the smooth base-change locus.  Shrinking a
constant-rank neighborhood gives
\begin{equation}\label{inv:horizontal-patch}
 \pi^*\theta\ge c g^*\omega_H,\qquad c>0,
\end{equation}
on a nonempty ordinary patch, disjoint from the Jacobian divisor.
Rank zero is impossible here when $\dim Y>0$.

On this equivariant model, with the trivial action on $Y$, the invariant
part of $g_*\mathcal O_Z(K_{Z/X})$ has rank one.  The natural Jacobian
section gives rank at least one.  If it were larger, global generation
after a sufficiently ample base twist would give two invariant
sections independent on the generic fiber.  Multiplying them by a
power of the invariant section of $b\pi^*B-g^*H$ puts them in
$K_{Z/X}+k b\pi^*B$.  Birational pushforward, using
$\pi_*\mathcal O_Z(K_{Z/X})=\mathcal O_X$, gives a ratio of
invariant sections of $kbB$ not in $F$, a contradiction.
This is invariant rank one; no assertion about the rank of the full
direct image has been made.

We have obtained the analytic data: smooth invariant coefficients,
invariant adjoint rank one, the natural nonvanishing Jacobian generator,
and the horizontal patch \eqref{inv:horizontal-patch}.  To obtain
unweighted integrability with exceptional support, we now use a
separate normal toroidal model.  Equivariance of that model is not
required because the meromorphic data on $X$ are already invariant.

\subsection{A klt model and its unweighted volume}
\label{inv:toroidal-volume}

By weak toroidalization and equidimensional subdivision
\cite{ADK2013,AK2000}, prepare
\[
 p:W\to X,\qquad f:W\to Y
\]
with $W$ normal and simplicial toroidal, $Y$ smooth, $f$
equidimensional and toroidal, and the $p$-exceptional locus in the
boundary.  The base may be modified birationally, and the smooth
equivariant graph resolution above is then replaced accordingly.
For clarity, the simplicial condition means that the primitive rays of each local
toric cone are linearly independent.  Such toric charts are finite
abelian quotients of smooth toric charts.  The simplicial refinement
can be made projectively using
generic heights on the existing rays.  Since the base cones are
simplicial and source rays map to base rays or zero, each subdivided
cone still maps onto a face of a base cone.  Equidimensionality is
preserved.  Simplicial toroidal charts have quotient singularities;
in particular $W$ is $\mathbb Q$-factorial and klt.

Put $D_j=\operatorname{div}(p^*\sigma_j)$ and
$V=(D_2-D_1)/(m_2-m_1)\sim_{\mathbb Q}p^*L$.
Equation~\eqref{inv:ratio-identity} gives
\[
 D_j=D_1+(m_j-m_1)V+
            \frac{f^*\operatorname{div}_Y(r_j)}{m_2-m_1}.
\]
Thus $V$ has nonnegative horizontal part.  Every vertical prime of $W$
dominates a base prime.  Its minimum divisor
\[
 M=\sum_Q\left(\min_{D\mapsto Q}
             \frac{\operatorname{ord}_D V}{\operatorname{ord}_D f^*Q}
        \right)Q
\]
is pseudoeffective by the coefficientwise finite-error argument
following \eqref{conv:affine-divisors}.  On $W$ one has the actual inequality
$V\ge f^*M$.

The smooth metric of $L$ defines a positive volume $\mu$ on $X$.
On a smooth base-change open, let $\rho$ be the coordinate density of
$g_*\pi^*\mu$.  This measure can equally be computed on $W_{\rm reg}$;
birational changes affect only sets of measure zero.  Its density is
the squared norm of the Jacobian adjoint generator divided by the base
canonical frame.  Compact averaging on the smooth direct image is a
holomorphic orthogonal projection onto the invariant line.  Its norm
is therefore the quotient norm.  For positive-dimensional fibers,
Lemma~\ref{coa:curvature} gives
\[
 i\partial\bar\partial(-\log\rho)\ge0,
\]
with strictness on a nonempty ordinary open by
\eqref{inv:horizontal-patch}.  If the relative dimension is zero,
connected general fibers make $g$ birational.  Over its isomorphism
open the integral metric is simply the transported coefficient metric;
its curvature gives the same assertion directly, including strictness
on the patch.  This use of a smooth coefficient and an orthogonal
quotient does not assert positivity of an arbitrary subbundle of a
singular direct image.

Write $E_W=K_{W/X}\ge0$.  For a base prime $Q$, set
\[
 a_D=\operatorname{ord}_D E_W,\qquad
 e_D=\operatorname{ord}_D f^*Q,\qquad
 b_Q=\max\left\{0,\min_{D\mapsto Q}
                         \frac{1+a_D}{e_D}-1\right\},
 \quad F_0=\sum_Q b_QQ.
\]
This effective rational divisor has finite support, and
$f^*F_0\le E_W$.  In fact, if $b_Q>0$, then
$e_Db_Q\le1+a_D-e_D\le a_D$ for every $D\mapsto Q$.
If a prime above $Q$ is nonexceptional over $X$, its $a_D$ is zero
and $b_Q=0$.  Hence $f^*F_0$ is exceptional over $X$.

On $-K_Y+F_0$ consider the weight
\[
 \Phi=-\log\rho+\sum_Qb_Q\log|q_Q|^2,
\]
initially off the indicated divisors, with $q_Q$ a local equation of
$Q$.  At a general point of a prime $D$ attaining the minimum, choose
coordinates with $f=(z_1^{e_D},z_2,\ldots,z_{\dim Y})$.
The pulled-back volume density is comparable to
$|z_1|^{2a_D}$ times a smooth positive density.  Integration of this
chart gives
\[
 \rho(y)\ge c|q_Q(y)|^{2((1+a_D)/e_D-1)}.
\]
Consequently $\Phi$ is locally bounded above at a general point of
every omitted divisor.  Plurisubharmonic extension, first across these
divisors and then across codimension two, gives a semipositive metric
on the rational line $-K_Y+F_0$, still smooth and strictly positive
on an ordinary open.

Its squared frame norm is unweighted locally integrable.  To check all
valuations, take a log resolution $h:W'\to W$ of the relevant data.
Integrating $e^{-\Phi}$ on the base is the same as integrating the
pulled-back volume divided by the local divisor factors of $f^*F_0$.
The exponent divisor upstairs is
\[
 K_{W'}-(ph)^*K_X-h^*f^*F_0
  =K_{W'/W}+h^*(E_W-f^*F_0).
\]
The first summand has coefficients greater than $-1$ because $W$ is
klt, and the second is effective because $E_W-f^*F_0$ is effective
and rationally Cartier.  On a normal-crossings resolution all
coefficients are therefore greater than $-1$, which is exactly the
local Jacobian integrability test.  This includes exceptional divisors
above higher-codimension centers.

Apply Corollary~\ref{conv:unweighted-effectivity} to the pseudoeffective
$M$.  It gives $F_0+\epsilon M\sim_{\mathbb Q}G\ge0$ for some
rational $\epsilon>0$.  Pullback and the inequality $V\ge f^*M$
make $V$ effective after a principal correction from $Y$ and a divisor
exceptional over $X$.  Equivalently apply Lemma~\ref{conv:return} on
a common resolution.  The underlying rational section is a power of
$\sigma_2/\sigma_1$ times a base function, so it is invariant.
Its extension on $X$ proves Theorem~\ref{inv:nonvanishing}.

\section{A maximal invariant base and signed adjoint metrics}
\label{inv:dominated-alternative}

There is another invariant construction when
$H^q(X,\mathcal O_X)=0$ for every $q>0$.  The maximal base is chosen
from all invariant rational maps whose polarizations are
pseudoeffectively dominated by $L$.  Its distinguishing output is a
sequence of signed boundaries, one for every positive adjoint twist.
The limit of those boundaries proves the pseudoeffectivity needed for
nonvanishing; the metric for a single twist would not suffice.

More precisely, if $X$ is smooth connected projective over $\mathbb C$
with $H^q(X,\mathcal O_X)=0$ for all $q>0$, a compact real torus
$T_c$ acts holomorphically, and the naturally linearized $L=-K_X$
is smoothly semipositive, this construction gives
$H^0(X,mL)^{T_c}\ne0$ for some $m>0$.  Since
$H^1(X,\mathcal O_X)=0$, the torus preserves the
isomorphism class of a very ample polarization and acts in the
projective linear stabilizer of its image.  A compact torus there lies
in an algebraic torus after conjugation; taking its Zariski closure
therefore permits equivariant projective graph resolutions.  Averaging
metric weights gives a smooth invariant coefficient.  The zero-index
value is one.  Proposition~\ref{inv:forms} and
Lemma~\ref{conv:determinant} supply a linearized $A$ with $-A$
pseudoeffective and sections as in
\eqref{inv:determinant-sequence}, unless a scalar invariant section
has already been obtained.

\subsection{Maximality and all positive adjoint ranks}

Consider invariant dominant rational maps to projective bases with
trivial action, subject on a resolution $\pi:Z\to X$, $g:Z\to Y$
to
\begin{equation}\label{inv:domination}
 \pi^*L-\delta g^*H\quad\hbox{pseudoeffective}
\end{equation}
for some ample $H$ and rational $\delta>0$.  Include the point map.
This condition is unchanged by replacing the resolution.  It persists
under birational modification of the base or its finite Stein
refinement: the pullback of an ample divisor under a generically finite
map is big and dominates a small ample class.  It also persists for the
joint image of two such maps, by addition and a product polarization.
Choose a map of maximal base dimension, take the relative algebraic
closure of its field in $\mathbb C(X)$, and resolve the base.  The
resulting field is still invariant, as in the finite-orbit argument
above.  Every other admissible invariant ratio belongs to it: its
adjoining cannot increase dimension, and the field is already
relatively algebraically closed.

Fix two exponents $a<b$ in the determinant sequence, set $m=b-a$, and
put $t=\sigma_b/\sigma_a$, an invariant rational section of $mL$.
For every larger exponent $i$, the ratio
\[
 \frac{\sigma_b^{i-a}}{\sigma_i^{b-a}\sigma_a^{i-b}}
\]
belongs to $\mathbb C(Y)$.  It is a pencil in a bundle $jL+dA$
with $j,d>0$, and $-A$ pseudoeffective makes that pencil satisfy
\eqref{inv:domination}.  Thus the horizontal orders of the
$\sigma_i$ vary affinely in $i$.  Their nonnegativity at unbounded
$i$ makes $D_t=\operatorname{div}(\pi^*t)$ effective horizontally.
If $Y$ is a point, this gives the desired section directly.

On a smooth equivariant graph resolution set $E=K_{Z/X}$.
For every integer $k\ge1$ the invariant generic adjoint rank is
\begin{equation}\label{inv:all-ranks}
 \operatorname{rank}\bigl(g_*\mathcal O_Z(E+km\pi^*L)\bigr)^{T_c}=1.
\end{equation}
The rational section $s_Et^k$ is regular on the generic fiber, hence
supplies a nonzero invariant element.  If the rank exceeded one,
global generation after twisting by $jg^*H$ would give two invariant
sections in $E+km\pi^*L+jg^*H$ whose ratio is not a base function.
Their effective divisors push to a pencil in
$D_0=kmL+j\pi_*g^*H$ on $X$; exceptionality removes $E$.
Pushing \eqref{inv:domination} forward gives an integer $b>0$ for
which $bL-D_0$ is pseudoeffective.  Resolve the pushed-down pencil by
$\sigma:U\to X$ and $h:U\to\mathbb P^1$.  Its fixed part $F$ is
effective and
$\sigma^*D_0\sim h^*\mathcal O_{\mathbb P^1}(1)+F$.  Consequently
\[
 b\sigma^*L-h^*\mathcal O_{\mathbb P^1}(1)
 \sim_{\mathbb Q}\sigma^*(bL-D_0)+F
 \quad\hbox{is pseudoeffective}.
\]
This is exactly the domination premise for the new invariant pencil.
Its ratio is not a base function, contradicting maximality.

We use a normal $\mathbb Q$-factorial simplicial toroidal
equidimensional model $p:W\to X$, $f:W\to Y$, as above, for the
divisorial calculations.  It need not carry the action.  All positivity
on invariant summands will instead be proved on the smooth equivariant
model, with base change and the ranks \eqref{inv:all-ranks}.

\subsection{The signed boundaries and their integrals}

For a prime $D\mapsto Q$ on $W$, put
\[
 d_D=\operatorname{ord}_D\operatorname{div}(p^*t),\qquad
 a_D=1+\operatorname{ord}_D K_{W/X},\qquad
 e_D=\operatorname{ord}_D f^*Q.
\]
Here $a_D\ge1$, and $a_D=1$ for nonexceptional primes.  Define
\[
 \gamma_Q=\min_{D\mapsto Q}\frac{d_D}{e_D},\qquad
 \lambda_Q(k)=\min_{D\mapsto Q}\frac{kd_D+a_D}{e_D}.
\]
The first minimum records the vertical slope of $t$; the second
includes the Jacobian order relevant to integration.  We choose the
coefficient $\beta_Q(k)$ of the base correction to satisfy
\[
 \lambda_Q(k)-1\le\beta_Q(k)<\lambda_Q(k),\qquad
 k\gamma_Q\le\beta_Q(k),\qquad
 \beta_Q(k)\le\frac{kd_D}{e_D}
        \quad(D\text{ nonexceptional}).
\]
These inequalities have different roles.  Near general points of
$D$ and $Q$, the local map $y_1=z_1^{e_D}$ turns a volume density
of order $2(kd_D+a_D-1)$ into a contribution to the base density
with exponent $2((kd_D+a_D)/e_D-1)$.  A prime attaining the minimum
therefore gives a lower bound with exponent $2(\lambda_Q(k)-1)$.
Dividing the base density by $|y_1|^{2\beta_Q(k)}$ adds
$\beta_Q(k)\log|y_1|^2$ to its negative logarithm.  The first
lower bound makes this corrected weight locally bounded above at $Q$.
The divided volume upstairs has exponent
$2(kd_D+a_D-1-e_D\beta_Q(k))$; the strict upper bound makes it
integrable at each prime over $Q$.  The toroidal calculation below
will check the remaining valuations.  The final upper bound ensures
that the section returned to $X$ has no pole at a nonexceptional
prime.  The comparison with $k\gamma_Q$ will be used after taking
the limit of the correction divisors.

All four requirements are met by
\[
 \beta_Q(k)=\max\{k\gamma_Q,\lambda_Q(k)-1\},\qquad
 B_k=\sum_Q\beta_Q(k)Q.
\]
The supports lie in one finite set.  Since every $a_D/e_D>0$,
both terms in the maximum are strictly below $\lambda_Q(k)$.
For a nonexceptional $D$, both $k\gamma_Q$ and
$\lambda_Q(k)-1\le kd_D/e_D+1/e_D-1$ are at most $kd_D/e_D$.
Finally, $\lambda_Q(k)-k\gamma_Q$ is bounded independently of $k$,
because the minimum is over a fixed finite set.  Thus
\begin{equation}\label{inv:beta-bounds}
 \beta_Q(k)<\lambda_Q(k),\qquad
 \beta_Q(k)\le\frac{kd_D}{e_D}\quad(D\text{ nonexceptional}),
 \qquad \frac{B_k}{k}\longrightarrow
 B_\infty:=\sum_Q\gamma_QQ.
\end{equation}
The chosen lower bound by $k\gamma_Q$ also gives $B_1\ge B_\infty$.
After the metrics show that $-K_Y+B_k$ is pseudoeffective for every
$k$, the displayed limit will imply pseudoeffectivity of $B_\infty$
and hence of $B_1$.

The pair
\begin{equation}\label{inv:sub-klt}
 (W,f^*B_k-k\operatorname{div}(p^*t)-K_{W/X})
\end{equation}
is sub-klt.  Horizontal coefficients are nonpositive.  A vertical
coefficient equals $e_D\beta_Q(k)-kd_D-(a_D-1)<1$.
If $e_D=1$, both terms in the maximum defining $\beta_Q(k)$ give
this coefficient at most zero.  Thus the positive part is supported on
the toroidal boundary.  On a simplicial toroidal chart the discrepancy
function has positive values $1-b_D$ on the boundary rays, and every
new ray has a positive linear combination of those values.  Hence an
effective boundary with coefficients below one is klt there.
Subtracting an effective rational divisor only improves discrepancies.
This proves \eqref{inv:sub-klt}, including all divisorial valuations.

Let $d\mu$ be the smooth anticanonical volume on $X$.  On the smooth
base-change locus the invariant rank-one generator of
$K_{Z/Y}+(km+1)\pi^*L$ is the relative form associated to
$s_Et^k$.  The squared norm in its coordinate frame is
\[
 \rho_k=\frac{g_*\pi^*(\|t\|^{2k}d\mu)}{dV_Y}.
\]
Smooth direct-image positivity applies to the smooth invariant
coefficient $(km+1)\pi^*L$.  Compact averaging is a holomorphic
orthogonal projection, so the invariant line has the quotient metric.
The squared frame norm on $G_k=-K_Y+B_k$ is
\[
 H_k=\rho_k\prod_Q|q_Q|^{-2\beta_Q(k)}.
\]
Near a prime attaining the minimum in $\lambda_Q(k)$, the local map
has form $y_1=z_1^{e_D}$ and the numerator volume is comparable to
$|z_1|^{2(kd_D+a_D-1)}$ times smooth volume.  Therefore
\[
 \rho_k(y)\ge c|q_Q(y)|^{2(\lambda_Q(k)-1)}.
\]
Since $\beta_Q(k)\ge\lambda_Q(k)-1$, the weight $-\log H_k$
is locally bounded above at general points of all missing divisors.
It extends plurisubharmonically there and across codimension two.
Thus every $G_k$ is pseudoeffective.

Moreover $H_k$ is unweighted locally integrable.  On a log resolution
$h:W'\to W$, the divisor of poles in the integrand is
\[
 h^*(f^*B_k-k\operatorname{div}(p^*t))-K_{W'/X}
 =h^*(f^*B_k-k\operatorname{div}(p^*t)-K_{W/X})-K_{W'/W}.
\]
Its coefficients are below one by \eqref{inv:sub-klt}.  The
normal-crossings Jacobian test and properness prove integrability
on the base.  This signed boundary calculation is separate from the
effective exceptional correction in Section~\ref{inv:toroidal-volume}.

\subsection{A global strict metric and the final section}

We give the strict-metric construction with its actual sheaf and ideal
conditions.  From \eqref{inv:domination} choose a singular metric on
$\pi^*L$ with curvature at least $\delta g^*\omega_H$.
Average its global weight difference $u$ from the smooth invariant
metric over $T_c$, in the sense of distributions.  This preserves the
quasiplurisubharmonic lower bound.  On $(m+1)\pi^*L$, add
$\epsilon u$ to the smooth coefficient weight.  For small
$\epsilon>0$ the resulting metric $h_\epsilon$ satisfies
\[
 \Theta_{h_\epsilon}\ge\epsilon\delta g^*\omega_H,
 \qquad\mathcal J(h_\epsilon)=\mathcal O_Z.
\]
Small-exponent integrability on compact $Z$ proves the ideal assertion.
Consequently the singular adjoint direct-image theorem
\cite[Theorem~3.3.5]{PT2018} acts on the full ordinary direct-image
sheaf.  Let $U\subset Y$ be a smooth base-change open on which this
sheaf is locally free.  Fubini makes the fiber integral norm finite
almost everywhere on $U$.

On $U$, averaging over the compact torus defines a holomorphic
projection $P$ onto the invariant summand.  Invariance of
$h_\epsilon$ makes $P$ orthogonal almost everywhere, where the integral
norm is finite.  The dual of the quotient embeds holomorphically in the
dual full bundle, and its dual norm is the restricted norm.  Logarithms
of these dual norms are plurisubharmonic by direct-image positivity.
They therefore define the semipositive quotient metric; its rank is
one by \eqref{inv:all-ranks}.  Subtracting a base potential before
applying this argument retains the lower curvature bound
$\epsilon\delta\omega_H$.  This proves positivity for this specified
quotient, without a claim about arbitrary singular invariant subbundles.

For $k=1$ the integral norm is obtained by inserting $e^{-\epsilon u}$
in the density $\rho_1$.  Since $u$ is locally bounded above, the
lower density estimates remain valid up to constants.  The same
extension across divisors, and then codimension two, gives a metric on
$G_1=-K_Y+B_1$ with a global strict curvature lower bound.
Together with the already constructed integrable semipositive metric,
strong openness and H\"older's inequality give a globally strict
metric on $G_1$ whose squared frame norms remain locally integrable.
Indeed take a sufficiently small positive fraction of the strict
weight; the integrable metric has some $L^{1+\eta}$ room by strong
openness, and the new singular weight has small-exponent integrability.

It remains to verify the algebraic premise for the adjoint class.
Every $G_k/k$ is pseudoeffective, and
$G_k/k=-K_Y/k+B_k/k\to B_\infty$.
The closedness of the pseudoeffective cone gives $B_\infty$
pseudoeffective.  Since $B_1\ge B_\infty$,
$B_1$ is pseudoeffective as well.  Choose a small ample rational $A'$
whose smooth curvature can be subtracted from the global strict lower
bound on $G_1$.  The resulting metric on $G_1-A'$ remains globally
strict and unweighted integrable.  Apply
Lemma~\ref{lem:weighted-resolution} with $J=0$ and a sufficiently
large divisible $r>1$.  On its log resolution, the fixed part divided
by $r$, minus the relative canonical divisor, has coefficients below
one.  A general member of the free part, divided by $r$, meets this
support transversely and has coefficient $1/r<1$.  It therefore gives
an effective klt $\Delta_0\sim_{\mathbb Q}G_1-A'$.
Add a general effective rational representative of $A'$ with small
coefficients, transverse on a common log resolution.  The resulting
$\Delta\sim_{\mathbb Q}G_1$ is effective, big, and klt.
Only the global strict bound and unweighted integrability are used
here; the mixed metric need not be smooth on an open set.
Now $K_Y+\Delta\sim_{\mathbb Q}B_1$ is pseudoeffective, so
big-boundary nonvanishing \cite[Theorem~D]{BCHM2010} gives
\[
 \operatorname{div}(h)+NB_1\ge0
\]
for some $N>0$ and $h\in\mathbb C(Y)^*$, after applying Lemma~\ref{lem:rational-recovery}.
The rational section $t^Nf^*h$ has no pole at horizontal primes.  At a
vertical prime nonexceptional over $X$ this follows from
$\beta_Q(1)\le d_D/e_D$ in \eqref{inv:beta-bounds}.
Equidimensionality handles every such prime on $W$.
It therefore extends to a nonzero section of $NmL$ on $X$.
Both $t$ and the base function are invariant, proving invariant
nonvanishing under the stated cohomology-vanishing hypothesis by this
second construction.

\section{Compact isometric monodromy and descent}
\label{inv:isometric-descent}

Let $X$ be a smooth connected projective complex variety whose
anticanonical bundle has a smooth Hermitian metric of semipositive
Chern curvature. The finite-cover structure theorem supplies a
compact factor to which the preceding invariant results apply.
There are two legitimate orders of operation: descend
invariant scalar sections, or descend twisted forms before applying an
ordinary conversion theorem.  We give the transport in both cases,
since the second does not assume an equivariant conclusion from
Theorem~\ref{conv:fixed-error}.

The finite-cover structure theorem
\cite[Theorem~2.1]{CompanionH}
supplies a
connected finite \'etale cover $\widehat X\to X$ whose universal cover
has the holomorphic isometric splitting
\[
 \widetilde X=\mathbb C^a\times R\times F.
\]
Here $R$ is a product of compact simply connected Ricci-flat factors,
$F$ is a product of the remaining compact projective irreducible
factors, and the deck group projects to a full translation lattice on
$\mathbb C^a$.  Its action on the compact factors lies in a compact
connected torus of holomorphic isometries, and the flat and Ricci-flat
factors have invariant nowhere-vanishing canonical frames.  Moreover
$H^q(F,\mathcal O_F)=0$ for $q>0$, and $L_F=-K_F$ is smoothly
semipositive.  Empty products and point factors are allowed.
These geometric and canonical-frame properties are supplied by the
cited theorem and its proof. The separate anticanonical nonvanishing
theorem of \cite[Theorem~1.1]{CompanionH} is not used here.

\begin{remark}\label{inv:tensor-holonomy}
One can strengthen the form-vanishing part of the structure description
to positive covariant tensors by a short local holonomy argument.
Let $S$ be one compact simply connected nonflat irreducible K\"ahler
factor, with nonnegative Ricci curvature, and let
$H\subset U(T_xS)$ be its holonomy group.  It is connected, and its
complex tangent representation is irreducible by de Rham
irreducibility.  If its determinant character is trivial, the canonical
connection is flat and $S$ is Ricci-flat; simple connectivity then
gives a parallel canonical frame.

Otherwise the differential of the determinant character is nonzero.
In the compact Lie algebra
$\mathfrak h=\mathfrak z\oplus[\mathfrak h,\mathfrak h]$, trace
vanishes on commutators.  Some central element therefore has nonzero
trace.  Schur's lemma makes that element a nonzero imaginary scalar,
whose exponential supplies the full scalar circle in $H$.
On $(T_x^*S)^{\otimes p}$ this circle has weight $-p$, so it has no
invariant vector for every $p>0$.  The K\"ahler Bochner formula with
nonnegative Ricci curvature makes every holomorphic covariant tensor
parallel.  Hence
\[
 H^0(S,(\Omega_S^1)^{\otimes p})=0\qquad(p>0).
\]
This argument concerns all covariant tensor powers, not only exterior
forms.  It is independent of the descent map below, which already
works for the alternating forms supplied by hard Lefschetz.
\end{remark}

For scalar descent, apply Theorem~\ref{inv:nonvanishing} to $F$,
or the compact-torus construction of
Section~\ref{inv:dominated-alternative}.  The compact action is
contained in an algebraic torus: $H^1(F,\mathcal O_F)=0$ preserves
a very ample polarization, and a compact abelian projective linear
group is conjugate into a diagonal torus.  Thus a positive power of
$L_F$ has a section fixed by the projected deck action.
Pull it to $\widetilde X$ and multiply its coefficient by the matching
inverse powers of the invariant canonical frames on $\mathbb C^a$
and $R$.  The resulting anticanonical section is deck-invariant and
nonzero, so it descends to $\widehat X$.  If $F$ is a point, begin
with its constant section.  The finite \'etale section norm then gives
a nonzero anticanonical multiple on $X$.

For the second order of operation, Proposition~\ref{inv:forms} gives a
fixed $p$ and invariant nonzero sections
\[
 u_m\in H^0(F,\Omega_F^p\otimes L_F^{m+1})
\]
at unbounded positive $m$, since the actual zero index equals one.
The pullback map for differential forms on the product is injective.
Tensoring the coefficient with the $(m+1)$-st inverse powers of the
canonical frames on the other factors gives a nonzero form with
coefficient $K_{\widetilde X}^{-(m+1)}$.  It is deck-invariant:
translations preserve the flat frame, compact monodromy preserves the
Ricci-flat frame, and $u_m$ is invariant on $F$.  It therefore descends
to a nonzero section of
\[
 \Omega_{\widehat X}^p\otimes(-(m+1)K_{\widehat X}).
\]
These holomorphic sections are algebraic because $\widehat X$ is
projective.  If $p=0$, we have a scalar section already.
Otherwise Corollary~\ref{conv:tensor-conversion} converts the descended
forms on the actual projective cover to an ordinary anticanonical
section.  The invariant information was used before this ordinary
conversion, not inferred from it afterward.

For clarity, the norm does not need a Galois cover.  If
$\nu:\widehat X\to X$ has degree $d$ and
$s\in H^0(\widehat X,-kK_{\widehat X})$ is nonzero, \'etaleness gives
$-kK_{\widehat X}=\nu^*(-kK_X)$.
Locally the product of $s$ over the $d$ sheets is invariant under their
permutation, so it glues to a section of $-dkK_X$.
At a general point every factor is nonzero, so the norm is nonzero.
This is \cite[Lemma~2.2]{CompanionH}, here with all
its hypotheses visible.

\subsection{Recovering an invariant factor section}

There is also a useful converse transport after ordinary conversion.
It identifies precisely the analytic role of isometric monodromy.

\begin{proposition}\label{inv:liouville-recovery}
Let $d>0$ be an integer.  Suppose a compact complex manifold $X'$ is the quotient of
$\mathbb C^a\times W\times F$ by a group acting by product maps,
whose translation projection is a full lattice in $\mathbb C^a$.
Assume $W,F$ are connected compact complex manifolds, the canonical
bundles of $\mathbb C^a$ and $W$ have invariant nowhere-vanishing
frames, and monodromy on $H^0(F,-dK_F)$ preserves a Hermitian norm.
A nonzero section of $-dK_{X'}$ determines a nonzero
monodromy-invariant section of $-dK_F$.
\end{proposition}
\begin{proof}
Lift the section and trivialize its coefficient on the first two
factors using the invariant frames.  It becomes a family of sections
of $-dK_F$.  This family is independent of $W$: its coefficients are
holomorphic functions on a connected compact manifold and hence
constant.  More explicitly, choose finitely many evaluations at points
of $F$, in local frames, giving an injective linear map from the
finite-dimensional space $H^0(F,-dK_F)$ to a coordinate space.  These
evaluations also show that the family is a holomorphic map
\[
 v:\mathbb C^a\longrightarrow H^0(F,-dK_F).
\]
Deck invariance and the preserved Hermitian norm make $\|v(z)\|$
periodic under the full translation lattice.  It is bounded on a
fundamental parallelepiped and hence everywhere.  Liouville's theorem
makes every coordinate of $v$ constant.  Its nonzero constant value is
fixed by every projected deck transformation.  The same proof includes
$a=0$.
\end{proof}
